% !tp-example Basic usage
\documentclass{article}
\begin{document}
\wcTpJavascriptViewer[src=/docs/components/_shared/intro-projects/javascript.json]{}
\end{document}

% !tp-example Javascript API
\documentclass{article}
\begin{document}
\wcScript[type=module]|
import "/components/javascript/javascript.js";
import "/components/javascript-viewer/javascript-viewer.js";

const example = document.createDocumentFragment();

const element1 = document.createElement("tp-javascript-viewer");
element1.setAttribute("src", "/docs/components/_shared/intro-projects/javascript.json");
example.append(element1);

document.body.append(example);
|
\end{document}

% !tp-example src
\documentclass{article}
\begin{document}
\wcTpJavascriptViewer[src=/docs/components/javascript-viewer/examples/example.json]{}
\end{document}

% !tp-example script
\documentclass{article}
\begin{document}
\wcTpJavascriptViewer{
  \wcScript[type=tp/javascript]|
        const app = document.querySelector('#app');

        if (app instanceof HTMLElement) {
          app.textContent = 'Hello JavaScript viewer';
        }

        console.log('Hello JavaScript');
      |
}
\end{document}

% !tp-example error
\documentclass{article}
\begin{document}
\wcTpJavascriptViewer{
  \wcScript[type=tp/javascript]|
        console.error(new Error('Intentional JavaScript error'));
      |
}
\end{document}

% !tp-example Single source file
\documentclass{article}
\begin{document}
\wcTpJavascriptViewer[src=/docs/components/_shared/single-source/example.js]{}
\end{document}

% !tp-example Compare presentation modes
\documentclass{article}
\begin{document}
\wcSection{
  \subsection*{mode="default" — Standard controls}
  Open Code to edit the source, then click Run or press \wcKbd*{Shift}+\wcKbd*{Enter} in the editor.

  Run to print a console message followed by an intentional error. In simple mode the console opens automatically. Comment out the message statements and run again to check that the error alone opens it; remove the error statement to try a successful run.

  \wcTpJavascriptViewer[mode=default]{
    \wcScript[type=tp/javascript,filename=main.js]|
document.body.innerHTML = "<p>The program ran.</p>";
console.log("Console message before the intentional error.");
throw new Error("Intentional error: the simple-mode console opens automatically.");
  |
  }
}

\wcSection{
  \subsection*{mode="lite" — Compact controls}
  Open Code to edit the source, then click Run or press \wcKbd*{Shift}+\wcKbd*{Enter} in the editor.

  Run to print a console message followed by an intentional error. In simple mode the console opens automatically. Comment out the message statements and run again to check that the error alone opens it; remove the error statement to try a successful run.

  \wcTpJavascriptViewer[mode=lite]{
    \wcScript[type=tp/javascript,filename=main.js]|
document.body.innerHTML = "<p>The program ran.</p>";
console.log("Console message before the intentional error.");
throw new Error("Intentional error: the simple-mode console opens automatically.");
  |
  }
}

\wcSection{
  \subsection*{mode="simple" — Editor and rendered output only}
  Edit the code below, keep the cursor in the editor and press \wcKbd*{Shift}+\wcKbd*{Enter} (Return) to run it. This viewer has no Run button.

  Run to print a console message followed by an intentional error. In simple mode the console opens automatically. Comment out the message statements and run again to check that the error alone opens it; remove the error statement to try a successful run.

  \wcTpJavascriptViewer[mode=simple]{
    \wcScript[type=tp/javascript,filename=main.js]|
document.body.innerHTML = "<p>The program ran.</p>";
console.log("Console message before the intentional error.");
throw new Error("Intentional error: the simple-mode console opens automatically.");
  |
  }
}

Programming viewers already have compact controls, so default and lite look alike. The first run may load the language engine.
\end{document}
